\documentclass[11pt]{article}
\usepackage[a4paper,margin=26mm]{geometry}
\usepackage[T1]{fontenc}
\usepackage{lmodern}
\usepackage{amsmath,amssymb,amsthm}
\usepackage{microtype}
\usepackage{array}
\usepackage{xurl}
\usepackage{hyperref}
\hypersetup{colorlinks=true,linkcolor=blue,urlcolor=blue,citecolor=blue,
 pdftitle={Lattice simplices whose Minkowski length exceeds the lattice diameter, and dilates whose Minkowski length is not quasi-linear},
 pdfauthor={Horace Dong}}
\newtheorem{theorem}{Theorem}
\newtheorem{lemma}{Lemma}
\newtheorem{proposition}{Proposition}
\newtheorem{corollary}{Corollary}
\theoremstyle{definition}
\newtheorem{definition}{Definition}
\theoremstyle{remark}
\newtheorem{remark}{Remark}
\newcommand{\repoversion}{1.25.0}
\setlength{\emergencystretch}{1.5em}
\newcommand{\Z}{\mathbb Z}
\newcommand{\R}{\mathbb R}
\newcommand{\Q}{\mathbb Q}
\newcommand{\N}{\mathbb N}
\DeclareMathOperator{\conv}{conv}
\DeclareMathOperator{\pyr}{pyr}

\title{Lattice simplices whose Minkowski length exceeds\\ the lattice diameter, and
dilates whose Minkowski length\\ is not quasi-linear}
\author{Horace Dong\\ \small Independent researcher\\
\small\texttt{maxwelldhx@gmail.com}\\
\small ORCID: \href{https://orcid.org/0009-0002-7554-7548}{0009-0002-7554-7548}}
\date{22 November 2026}
\begin{document}
\maketitle

\begin{abstract}
The Minkowski length $L(P)$ of a lattice polytope $P$ is the largest number of
lattice polytopes of positive dimension whose Minkowski sum is contained in
$P$; the lattice diameter $\ell(P)$ is one less than the largest number of
collinear lattice points in $P$. Soprunov and Soprunova proved that
$L(tP)$ is eventually quasi-linear in $t$, and asked whether $L(tP)$ is always
quasi-linear (their Problem~2) and whether $L=\ell$ for every simplex
(their Problem~3). We answer both questions in the negative. The lattice
tetrahedron with vertices $(1,0,1)$, $(3,3,1)$, $(2,0,3)$, $(0,3,3)$ has
$\ell=1$ and $L=2$. The empty lattice tetrahedron $U$ with vertices $0$,
$(1,0,0)$, $(0,0,1)$, $(13,18,1)$ has $L(U)=1$ and $L(tU)=3t$ for every
integer $t\ge2$, so $L(tU)$ is not quasi-linear. Lattice pyramids carry both
examples to every dimension $d\ge3$. Since equality $L=\ell$ is known for
lattice triangles, the answer to Problem~3 is: true for $d\le2$, false for
every $d\ge3$. The results for $d\ge3$ are formally verified in Lean~4. The
proofs were found and formalized by AI agents working under the author's
direction; their role is described at the end of the paper.
\end{abstract}

\section{Introduction}

We work in $\R^d$ with the integer lattice $\Z^d$. A \emph{lattice polytope}
is the convex hull of finitely many points of $\Z^d$; it has positive
dimension if it contains two distinct points. For a lattice polytope
$P\subset\R^d$, Soprunov and Soprunova~\cite[Definition~1.1]{ss} define the
\emph{Minkowski length} $L(P)$ as ``the largest number of lattice polytopes
$Q_1,\dots,Q_L$ of positive dimension whose Minkowski sum is contained in
$P$''. The \emph{lattice diameter} $\ell(P)$ is one less than the largest
number of collinear lattice points in $P$~\cite[Introduction and p.~5]{ss}.
Always $\ell(P)\le L(P)$, since $\ell(P)$ is the case of collinear summands
($\ell(P)=L_1(P)$ in the notation of~\cite[Definition~2.5]{ss}). For an
integer $t\ge1$, $tP=\{tx\mid x\in P\}$.

Following~\cite[Definition~1.3]{ss}, a function $f\colon\N\to\Q$ is
\emph{quasi-linear} if there are $k\in\N$ and linear polynomials
$p_0,\dots,p_{k-1}\in\Q[t]$ such that $f(t)=p_r(t)$ whenever $t\equiv r\bmod
k$, and \emph{eventually quasi-linear} if $f$ agrees with a quasi-linear
function for all large enough $t$. The main theorem of~\cite{ss}
(Theorem~2.20) states that $t\mapsto L(tP)$ is eventually quasi-linear for
every lattice polytope $P$. The introduction of~\cite{ss} describes this
result as ``eventually quasi-polynomial with linear constituents (we say
`quasi-linear' for short)'', keeping the word `eventually'; Definition~1.3 and
the discussion before Problem~2 distinguish the two notions, and the authors
note that a certain identity ``would imply that $L(tP)$ is not just
eventually quasi-linear, but quasi-linear''~\cite[p.~12]{ss}. In Section~4
of~\cite{ss} they ask:

\begin{quote}
\textbf{Problem 2 of~\cite{ss}.} Prove that $L(tP)$ is quasi-linear or give an
example of a lattice polytope $P$ for which $L(tP)$ is not quasi-linear.

\textbf{Problem 3 of~\cite{ss}.} Prove or disprove that for any simplex in
$\R^d$ its Minkowski length coincides with its lattice diameter.
\end{quote}

\noindent Here ``simplex'' means a lattice simplex, since $L$ is defined for
lattice polytopes. Numbers and page numbers of items in~\cite{ss} refer to
its arXiv version (arXiv:1412.4404v2, 23 May 2016), where Problems~1--3 are
on p.~12. We checked
that Problems~1--3 have the same numbers and wording in arXiv:1412.4404v1. A
Codex agent reported checking the numbers and wording against Section~4 of
the journal version; we have not seen the publisher's text ourselves.

Problem~3 was motivated by two cases in which equality holds: the dilates
$t\Delta$ of a unimodular simplex, for which $L(t\Delta)=\ell(t\Delta)=t$
\cite[Theorem~2.1, Example~2.6]{ss}, and lattice triangles~\cite[Corollary~3.2]{ss}.
We show that it fails in every dimension $d\ge3$, and that Problem~2 has a
negative answer.

\begin{theorem}\label{thm:diam}
Let $T=\conv\{(1,0,1),(3,3,1),(2,0,3),(0,3,3)\}\subset\R^3$. Then $T$ is a
$3$-dimensional lattice simplex with $\ell(T)=1$ and $L(T)=2$. More generally,
for every $d\ge3$ there is a $d$-dimensional lattice simplex $S\subset\R^d$
with $\ell(S)=1$ and $L(S)=2$.
\end{theorem}

\begin{theorem}\label{thm:dil}
Let $U=\conv\{(0,0,0),(1,0,0),(0,0,1),(13,18,1)\}\subset\R^3$. Then $U$ is a
$3$-dimensional lattice simplex with $L(U)=1$ and $L(tU)=3t$ for every integer
$t\ge2$; in particular $t\mapsto L(tU)$ ($t\ge1$) is not quasi-linear. More
generally, for every $d\ge3$ there is a $d$-dimensional lattice simplex
$S\subset\R^d$ with $L(S)=1$ and $L(tS)=3t$ for all integers $t\ge2$, so that
$L(tS)$ is not quasi-linear.
\end{theorem}

Both theorems are formally verified in Lean~4 (Section~\ref{sec:lean}). The
sequence $L(tU)$ is linear from $t=2$ on, in agreement with the eventual
quasi-linearity proved in~\cite{ss} (Remark~\ref{rem:period}); only its first
term breaks quasi-linearity. Together with the known cases $d\le2$ we obtain
the following answer to Problem~3. Only the part $d\ge3$ is in the Lean files.

\begin{corollary}\label{cor:p3}
Every lattice simplex $S\subset\R^d$ with $d\le2$ satisfies $L(S)=\ell(S)$:
for $d=1$ this is trivial, and for $d=2$ it is~\cite[Corollary~3.2]{ss}. For
every $d\ge3$ there is a $d$-dimensional lattice simplex with $L(S)\neq\ell(S)$.
\end{corollary}

\subsection*{Prior work}
The Minkowski length was introduced by Soprunov and Soprunova in connection
with toric surface codes (see~\cite{ss} and the references there); in
dimension three it can be computed by the algorithm of Beckwith, Grimm,
Soprunova and Weaver~\cite{bgsw}. Both tetrahedra above are known objects.
The tetrahedron $T$ has $8$ lattice points and lattice width~$2$
(Section~\ref{sec:T}), so its lattice equivalence class is one of the $193$
tetrahedra among the $2675$ lattice $3$-polytopes of width larger than one
with $8$ lattice points enumerated by Blanco and Santos~\cite[Theorem~4 and
Table~2]{bs} (more precisely, since $T$ has four interior lattice points, one
of the $90$ such tetrahedra with four interior lattice points,
\cite[Table~5]{bs}). Their Section~5.6 treats the polytopes of Minkowski length one
(the ``dps'' polytopes); $T$ has Minkowski length two and is not among them.
We have not identified the entry of $T$ in their database. The tetrahedron $U$
is empty, that is, its only lattice points are its vertices
(Lemma~\ref{lem:Upoints}); empty lattice tetrahedra were classified by
White~\cite{white}, and $U$ has normalized volume~$18$. The closest later work
we know is by Meyer, Soprunov and Soprunova~\cite{mss}, who study maximal
decompositions in $3$-polytopes and, in particular, $L(P+Q)$ for empty
tetrahedra $P,Q$ \cite[Lemma~3.2 and Theorem~3.4]{mss}. Their Theorem~3.4,
with $P=Q$, together with superadditivity, gives $L(2P)\ge3$ for every empty
tetrahedron $P$ not equivalent to one of their two tetrahedra of normalized
volume $1$ and~$2$, hence for $U$. It says nothing about the whole sequence
$L(tU)$ or about quasi-linearity, and~\cite{mss} does not discuss Problems~2
or~3 of~\cite{ss}. As far as we know, Problems~2 and~3 have not been answered
elsewhere; the searches are described in the final section.

\section{Segments and corners}\label{sec:prelim}

As noted in~\cite[p.~3]{ss}, each summand may be shrunk to a lattice segment.
We use the following form of this remark. For lattice points
$a_1,\dots,a_k,b_1,\dots,b_k$ and $\varepsilon\in\{0,1\}^k$, call
$\sum_i c_i$ with $c_i=a_i$ if $\varepsilon_i=0$ and $c_i=b_i$ if
$\varepsilon_i=1$ a \emph{corner}.

\begin{lemma}\label{lem:corners}
Let $P\subseteq\R^d$ be convex and $k\ge1$. Then $P$ contains a Minkowski sum
of $k$ lattice polytopes of positive dimension if and only if there are lattice
points $a_i\ne b_i$ $(1\le i\le k)$ all of whose $2^k$ corners lie in $P$.
The same holds if the corners are only known to lie in $P$ after adding a fixed
lattice vector $g$.
\end{lemma}

\begin{proof}
If $Q_1+\dots+Q_k\subseteq P$, choose two distinct lattice points $a_i,b_i$ of
$Q_i$; every corner lies in the sum. Conversely, the sum
$[a_1,b_1]+\dots+[a_k,b_k]$ is the convex hull of the corners, hence contained
in $P$. A lattice vector $g$ is absorbed by replacing $a_1,b_1$ with
$a_1+g,b_1+g$.
\end{proof}

Thus $L(P)$ is the largest $k$ for which such points exist; in particular a
lattice translation of the whole sum is allowed, but not a translation by a
non-lattice vector.

\section{The tetrahedron \texorpdfstring{$T$}{T}}\label{sec:T}

\begin{lemma}\label{lem:Tpoints}
For $p=(x,y,z)\in\R^3$ put
\[
\beta(p)=\bigl(39-6x-4y-9z,\ 6x+4y-3z-3,\ 6x-4y+9z-15,\ -6x+4y+3z+3\bigr).
\]
Then $\beta(p)/24$ is the vector of barycentric coordinates of $p$ with respect
to the vertices of $T$, in the order listed, so $T=\{p:\beta(p)\ge0\}$. The
determinant of the edge vectors at $(1,0,1)$ is $-24$, so $T$ is
$3$-dimensional. The lattice points of $T$ are its four vertices and
$(1,1,2)$, $(1,2,2)$, $(2,1,2)$, $(2,2,2)$.
\end{lemma}

\begin{proof}
The four components are affine, sum to $24$, and take the value $24$ at the
corresponding vertex and $0$ at the other three. From $\beta\ge0$,
$\beta_0+\beta_1=36-12z$, $\beta_2+\beta_3=12z-12$, $\beta_1+\beta_3=8y$ and
$\beta_0+\beta_2=24-8y$ give $1\le z\le3$ and $0\le y\le3$. For $z=1$,
$\beta_2,\beta_3\ge0$ force $6x=4y+6$, so $(x,y)\in\{(1,0),(3,3)\}$; for $z=3$,
$\beta_0,\beta_1\ge0$ force $6x+4y=12$, so $(x,y)\in\{(2,0),(0,3)\}$; for $z=2$
the conditions $9\le6x+4y\le21$ and $-3\le6x-4y\le9$ leave exactly
$x,y\in\{1,2\}$.
\end{proof}

The four points at height $2$ form a unit square; all components of $\beta$
are positive at them, so they are interior points of $T$. Hence
every primitive linear functional takes at least three values on $T\cap\Z^3$,
and the lattice width of $T$ is $2$ (attained by $z$).

\begin{proposition}\label{prop:T}
$\ell(T)=1$ and $L(T)=2$.
\end{proposition}

\begin{proof}
\emph{$\ell(T)=1$.} The eight lattice points of $T$ lie in eight different
classes modulo~$2$. If three lattice points of $T$ were collinear, let $p$ and
$q$ be the outer two and $v$ the primitive vector in the direction $q-p$; then
$q=p+mv$ with $m\ge2$, and $p$ and $p+2v$ are distinct lattice points of $T$
(by convexity) that are congruent modulo~$2$, a contradiction. Two vertices
give two distinct lattice points.

\emph{$L(T)\ge2$.} The square $(1,1,2)+[0,e_1]+[0,e_2]$ lies in $T$.

\emph{$L(T)\le2$.} Suppose $T$ contains a sum of three or more summands. By
Lemma~\ref{lem:corners}, keeping three of the segments and fixing an endpoint
of each of the others, there are a lattice point $p$ and nonzero lattice
vectors $u_1,u_2,u_3$ such that $p+\sum_{i\in I}u_i\in T$ for all
$I\subseteq\{1,2,3\}$. Choose a linear functional $\varphi$ with
$\varphi(u_i)\ne0$ for all $i$; replacing $p$ by $p+u_i$ and $u_i$ by $-u_i$
does not change the set of these eight points, so we may assume
$\varphi(u_i)>0$. If $u_1,u_2,u_3$ are linearly dependent modulo~$2$, some
nonempty $I$ has $\sum_{i\in I}u_i\equiv0\bmod 2$; this sum is nonzero since
its value under $\varphi$ is positive, so $p$ and $p+\sum_{i\in I}u_i$ are
distinct congruent lattice points of $T$, which is impossible. Otherwise the
eight points are pairwise incongruent modulo~$2$, hence they are exactly the
eight lattice points of $T$, including its vertices. The zonotope
$Z=p+[0,u_1]+[0,u_2]+[0,u_3]$ is then contained in $T$ and contains its
vertices, so $Z=T$. But $Z$ is centrally symmetric and $T$ is not: a central
symmetry of $T$ would permute the vertices without fixed points (the centre is
an interior point), so the four vertices would form two pairs with a common
midpoint, that is, a parallelogram, and would be coplanar.
\end{proof}

The Lean proof of $L(T)\le2$ uses a different, finite argument. Let
$\kappa(x,y,z)=x+32y+1024z$. On the eight lattice points of $T$, $\kappa$
takes the values $\kappa_1,\dots,\kappa_8$ given by
\[
\begin{array}{@{}c|cccccccc@{}}
\text{point} & (1,0,1)&(3,3,1)&(2,0,3)&(0,3,3)&(1,1,2)&(1,2,2)&(2,1,2)&(2,2,2)\\
\kappa & 1025&1123&3074&3168&2081&2113&2082&2114
\end{array}
\]
With $p,u_i$ as above, let
$p,p+u_1,p+u_2,p+u_3$ be the points number $i,j,k,l$ of this list; since
$u_1,u_2,u_3\ne0$, $i$ differs from $j$, $k$ and $l$. By linearity, the values
of $\kappa$ at the other four points $p+u_1+u_2,\dots,p+u_1+u_2+u_3$ are
$\kappa_j+\kappa_k-\kappa_i$, $\kappa_j+\kappa_l-\kappa_i$,
$\kappa_k+\kappa_l-\kappa_i$ and $\kappa_j+\kappa_k+\kappa_l-2\kappa_i$, so
these four numbers must occur in the list. A computation over all $8^4$
quadruples $(i,j,k,l)$, carried out by the Lean kernel, shows that this happens
only when $i\in\{j,k,l\}$. Only the linearity of $\kappa$
is used, not its injectivity on~$\Z^3$; the $u_i$ need not be distinct.

\section{The tetrahedron \texorpdfstring{$U$}{U}}\label{sec:U}

For $p=(x,y,z)\in\R^3$ put
\[
F(p)=\bigl(-18z-18x+13y,\ 18x-13y,\ 18z-y,\ y\bigr).
\]
The components of $F$ are linear and sum to~$0$. At the vertices $0$, $e_1$,
$e_3$, $(13,18,1)$ of $U$, $F$ takes the values $(0,0,0,0)$, $(-18,18,0,0)$,
$(-18,0,18,0)$, $(-18,0,0,18)$, so for an integer $t\ge1$ the numbers
$\bigl(18t+F_0(p),F_1(p),F_2(p),F_3(p)\bigr)/(18t)$ are the barycentric
coordinates of $p$ with respect to the vertices of $tU$, and
\begin{equation}\label{eq:tU}
tU=\{p\in\R^3 : F_0(p)\ge-18t,\ F_1(p)\ge0,\ F_2(p)\ge0,\ F_3(p)\ge0\}.
\end{equation}
The determinant of the edge vectors at $0$ is $-18$, so $U$ is
$3$-dimensional.

\begin{lemma}\label{lem:Upoints}
The only lattice points of $U$ are its four vertices, and $L(U)=1$.
\end{lemma}

\begin{proof}
By~\eqref{eq:tU} with $t=1$, $F_0+F_1=-18z\ge-18$ and $F_2+F_3=18z\ge0$ give
$z\in\{0,1\}$, and $0\le y\le18z$. For $z=0$ we get $y=0$ and $0\le18x\le18$;
for $z=1$ we get $13y\le18x\le13y$, so $18x=13y$ and $y\in\{0,18\}$. This
leaves the four vertices. The edge $[0,e_1]$ gives $L(U)\ge1$. If $U$
contained a sum of two summands, Lemma~\ref{lem:corners} would give lattice
points $a,b,c,d$ of $U$ with $a+d=b+c$, $a\ne b$ and $a\ne c$. Among four
affinely independent points this is impossible: it would give either three
distinct collinear points ($b=c$ or $a=d$) or four coplanar ones.
\end{proof}

For a lattice vector $u\ne0$ let
\[
\operatorname{cost}(u)=\sum_{j=0}^{3}\max\bigl(0,-F_j(u)\bigr)
=\tfrac12\sum_{j=0}^{3}|F_j(u)| .
\]

\begin{lemma}\label{lem:cost}
$\operatorname{cost}(u)\ge6$ for every $u\in\Z^3\setminus\{0\}$.
\end{lemma}

\begin{proof}
Let $u=(x,y,z)$. If $z\ne0$, then
$\sum_j|F_j(u)|\ge|F_0(u)+F_1(u)|+|F_2(u)+F_3(u)|=36|z|$, so
$\operatorname{cost}(u)\ge18$. If $z=0$, then
$\operatorname{cost}(u)=|18x-13y|+|y|$. This is at least $18$ if $y=0$ (then
$x\ne0$) and at least $6$ if $|y|\ge6$. For $|y|=1,2,3,4,5$ the distance from
$13|y|$ to the nearest multiple of $18$ is $5,8,3,2,7$, which gives
$\operatorname{cost}(u)\ge6,10,6,6,12$.
\end{proof}

\begin{proposition}\label{prop:U}
$L(tU)=3t$ for every integer $t\ge2$, and $t\mapsto L(tU)$ $(t\ge1)$ is not
quasi-linear.
\end{proposition}

\begin{proof}
\emph{Upper bound, for every $t\ge1$.} Let $[a_i,b_i]$ $(1\le i\le k)$ be
lattice segments, $a_i\ne b_i$, all of whose corners lie in $tU$. Fix $j$ and
choose for each $i$ the endpoint at which $F_j$ is smaller; the corresponding
corner lies in $tU$, so by~\eqref{eq:tU}
\[
\sum_{i=1}^k\min\bigl(F_j(a_i),F_j(b_i)\bigr)\ge-18t\ \text{ for } j=0,
\qquad \ge0\ \text{ for } j=1,2,3 .
\]
Add the four inequalities and use $\sum_jF_j(a_i)=0$:
$\sum_i\sum_j\bigl(F_j(a_i)-\min(F_j(a_i),F_j(b_i))\bigr)\le18t$. The inner
sum is $\operatorname{cost}(b_i-a_i)\ge6$ by Lemma~\ref{lem:cost}, so
$6k\le18t$, i.e.\ $k\le3t$. The argument allows arbitrary lattice translations
and directions.

\emph{Lower bound, for every $t\ge2$.} By~\eqref{eq:tU}, the points $(1,0,1)$
and $(13,18,1)=(1,0,1)+6\,(2,3,0)$ lie in $2U$, and $(0,0,2)$ and
$(27,36,2)=(0,0,2)+9\,(3,4,0)$ lie in $3U$. So $2U$ contains a sum of $6$
lattice segments and $3U$ a sum of $9$. Every $t\ge2$ can be written as
$t=2i+3j$ with $i,j\ge0$, and $tU=i\cdot2U+j\cdot3U$ by convexity, so $tU$
contains a sum of $6i+9j=3t$ lattice segments.

\emph{Not quasi-linear.} Suppose $L(tU)=\alpha_rt+\gamma_r$ for $t\equiv r\bmod k$,
with $k\ge1$ and rational $\alpha_r,\gamma_r$. The integers $1$, $1+k$ and
$1+2k$ lie in one residue class. Since $L(tU)=3t$ at $t=1+k$ and $t=1+2k$,
that class has $\alpha=3$ and $\gamma=0$, so $L(U)=3$, contradicting
Lemma~\ref{lem:Upoints}.
\end{proof}

\begin{remark}\label{rem:period}
The sequence agrees with~\cite[Theorem~2.20]{ss}. By
Proposition~\ref{prop:U}, the rational Minkowski length
$\lambda(U)=\sup_tL(tU)/t$ \cite[Definition~2.7]{ss} equals $3$, and the
period of $U$ in the sense of~\cite[Definition~2.17]{ss} is $k=2$. Theorem~2.20
of~\cite{ss} then gives $L(tU)=6\lfloor t/2\rfloor+c_r$ for large
$t\equiv r\bmod2$, with $L(rU)\le c_r\le 3r$; here $c_0=0$ and $c_1=3$, while
$L(U)=1$. This remark is not stated in the Lean files.
\end{remark}

\section{Lattice pyramids}\label{sec:pyr}

Let $P=\conv\{v_0,\dots,v_d\}\subset\R^d$ ($d\ge1$) be a $d$-dimensional
lattice simplex. Write points of $\R^{d+1}$ as $(h,y)$ with $h\in\R$ and
$y\in\R^d$, and let
\[
\pyr P=\conv\{(1,0),(0,v_0),\dots,(0,v_d)\}\subset\R^{d+1},
\]
the lattice pyramid of height one over $P$ with apex $(1,0)$.

\begin{lemma}\label{lem:pyr}
(a) $\pyr P$ is a $(d+1)$-dimensional lattice simplex.
(b) For every integer $t\ge1$,
$t\pyr P=\{(h,y): 0\le h\le t,\ y\in(t-h)P\}$, where $0P=\{0\}$.
(c) The lattice points of $\pyr P$ are the apex $(1,0)$ and the points $(0,q)$
with $q\in P\cap\Z^d$.
\end{lemma}

\begin{proof}
(a) An affine dependence among the vertices has zero apex weight (look at the
coordinate $h$), and then it is an affine dependence among the $v_i$. (b) A
point of $t\pyr P$ is a combination of the vertices with nonnegative weights
of total mass $t$; the apex weight is $h$, and the base weights have total
mass $t-h$ and combine to $y$. (c) follows from (b) with $t=1$, since a
lattice point has $h\in\{0,1\}$.
\end{proof}

\begin{proposition}\label{prop:pyrdiam}
If $\ell(P)=1$, then $\ell(\pyr P)=1$.
\end{proposition}

\begin{proof}
Two distinct lattice points of $P$ give two of $\pyr P$. Let three lattice
points of $\pyr P$ be collinear. If all lie at height $0$, they come from three
collinear lattice points of $P$, two of which coincide. Otherwise one of them
is the apex; a line through the apex meets the hyperplane $h=0$ in at most one
point, so the other two coincide.
\end{proof}

\begin{proposition}\label{prop:pyrlen}
For every integer $t\ge1$, $L(t\pyr P)=L(tP)$.
\end{proposition}

\begin{proof}
If segments $[a_i,b_i]$ have all corners in $tP$, the segments
$[(0,a_i),(0,b_i)]$ have all corners in $t\pyr P$ by Lemma~\ref{lem:pyr}(b).
Conversely, let $[a_i,b_i]$ $(1\le i\le k)$ be lattice segments in
$\Z^{d+1}$, $a_i\ne b_i$, with all corners in $t\pyr P$, labelled so that
$h(b_i)\ge h(a_i)$. Call a segment \emph{rising} if $h(b_i)>h(a_i)$, and let
$m$ be the number of rising segments. All corners in which every rising
segment takes its upper endpoint $b_i$ have the same height $H$, because the
other segments are horizontal. The lowest corner has height $\ge0$ and each
rising segment adds at least~$1$, so $m\le H\le t$. By
Lemma~\ref{lem:pyr}(b), for every choice of endpoints on the horizontal
segments, the $y$-parts of these corners lie in $(t-H)P$; they are
$g+\sum_{i\text{ horizontal}}y(c_i)$, where $g$ is the sum of the $y$-parts
of the upper endpoints of the rising segments. Now replace each rising
segment by $[v_0,v_d]$ and add the lattice vector $g+(H-m)v_0$. Every corner
of the new configuration of $k$ segments lies in $(t-H)P+HP=tP$, by convexity
of $P$. By Lemma~\ref{lem:corners}, $L(tP)\ge k$.
\end{proof}

\begin{proof}[Proof of Theorems~\ref{thm:diam} and~\ref{thm:dil}]
The statements about $T$ and $U$ are Lemma~\ref{lem:Tpoints},
Proposition~\ref{prop:T}, Lemma~\ref{lem:Upoints} and
Proposition~\ref{prop:U}. Put $S_3=T$ (respectively $S_3=U$) and
$S_{d+1}=\pyr S_d$. By Lemma~\ref{lem:pyr}(a), $S_d$ is a $d$-dimensional
lattice simplex; by Proposition~\ref{prop:pyrlen}, $L(tS_d)=L(tS_3)$ for every
$t\ge1$, and by Proposition~\ref{prop:pyrdiam}, $\ell(S_d)=\ell(S_3)$ when
$S_3=T$. The argument that $L(tU)$ is not quasi-linear uses only $L(U)=1$
and $L(tU)=3t$ for $t\ge2$, so it applies to $S_d$.
\end{proof}

\section{Remarks}\label{sec:remarks}

\paragraph{The identity \texorpdfstring{$L_1(P)=\lfloor\lambda_1(P)\rfloor$}{L1(P)=floor(lambda1(P))}.}
In Definition~2.10 the source introduces the rational diameter $s(P)$,
defined through segments $[0,sv]$ contained in $P$ ``after a translation by a
rational vector''; after it, the source shows $s(P)=\lambda_1(P)$ and
concludes
$L_1(P)=\lfloor\lambda_1(P)\rfloor$ from $\ell(P)=\lfloor s(P)\rfloor$
\cite[p.~6]{ss} (the same text is in v1). The simplex $U$ shows that this
identity does not hold in general. By Lemma~\ref{lem:Upoints},
$L_1(U)=\ell(U)=1$. On the other hand, the seven lattice points
$(1,0,1)+j(2,3,0)$, $0\le j\le6$, lie in $2U$, so $L_1(2U)\ge6$ and
$\lambda_1(U)\ge3$; with Proposition~\ref{prop:U}, $\lambda_1(U)=3$. The
segment $\frac12(1,0,1)+[0,3(2,3,0)]$ lies in $U$ but has non-integral
endpoints. The cause is that $s(P)$ allows a translation by a rational vector,
whereas $\ell(P)$ counts lattice points, that is, allows only lattice
translations of lattice segments; the step $\ell(P)=\lfloor s(P)\rfloor$
fails for $U$. Remark~2.21 of~\cite{ss} uses this identity to conclude that
``the function $L_1(tP)$ is, in fact, quasi-linear''. That conclusion is
therefore not supported in general. (This paper does not formalize the
values of $L_1(tU)$.) The main
theorem~\cite[Theorem~2.20]{ss} is not affected; its proof does not use the
identity. Nor is \cite[Corollary~3.2]{ss}, used in Corollary~\ref{cor:p3}: in
a triangle a longest chord in a given direction can be chosen with a vertex as
an endpoint, so the segment used there is a lattice translate of
$[0,\lfloor s(T)\rfloor v]$. This paragraph is not part of the Lean files.

\paragraph{What remains open.}
Problem~1 of~\cite{ss}, on the supremum of $\lambda(P)-L(P)$, is not
addressed here. Our answer to Problem~2 is a $3$-dimensional example and its
pyramids. In dimension one $L(tP)=tL(P)$ is linear. For lattice polygons the
authors report that $L(tP)$ was quasi-linear in all examples they
computed~\cite[p.~12]{ss}; whether this holds for every lattice polygon is,
as far as we know, still open.

{\sloppy
\section{Formal verification}\label{sec:lean}

Theorems~\ref{thm:diam} and~\ref{thm:dil} are formally verified in Lean~4
(version 4.34.1) with Mathlib~\cite{mathlib} (commit \texttt{d13f23b7}). The
$16$ files are in version~\repoversion{} of the repository~\cite{repo},
directory \texttt{lean/Research/TaskP/GeomSimplex/}. Corollary~\ref{cor:p3}
for $d\le2$, Remark~\ref{rem:period} and Section~\ref{sec:remarks} are not
formalized.

\paragraph{The statements.}
The statement file \texttt{Challenge.lean} contains only definitions and
statements, no proofs. In ASCII transliteration (the docstrings and the
\texttt{import}, \texttt{namespace} and \texttt{end} lines omitted; the file
itself uses Unicode, e.g.\ $\mathbb Z$, $\mathbb R$, $\sum$, $\to$ and $\ne$):
{\footnotesize
\begin{verbatim}
abbrev Lattice (d : Nat) := Fin d -> Int
abbrev Space (d : Nat) := Fin d -> Real
def realPoint {d : Nat} (a : Lattice d) : Space d := fun i => a i
def polytope {d : Nat} (vertices : Finset (Lattice d)) : Set (Space d) :=
  convexHull Real (realPoint '' (vertices : Set (Lattice d)))
def simplex {d : Nat} (vertices : Fin (d + 1) -> Lattice d) : Set (Space d) :=
  convexHull Real (Set.range (realPoint o vertices))
def fullSimplex {d : Nat} (vertices : Fin (d + 1) -> Lattice d) : Prop :=
  AffineIndependent Real (realPoint o vertices)
def dilate {d : Nat} (t : Nat) (P : Set (Space d)) : Set (Space d) :=
  (fun x => (t : Real) smul x) '' P
def packs {d : Nat} (P : Set (Space d)) (k : Nat) : Prop :=
  exists V : Fin k -> Finset (Lattice d),
    (forall i, 1 < (V i).card) /\
    forall x : Fin k -> Space d, (forall i, x i in polytope (V i)) -> sum i, x i in P
def lengthIs {d : Nat} (P : Set (Space d)) (n : Nat) : Prop :=
  packs P n /\ forall k, packs P k -> k <= n
def diameterOne {d : Nat} (P : Set (Space d)) : Prop :=
  (exists a b : Lattice d, a != b /\ realPoint a in P /\ realPoint b in P) /\
  forall a b c : Lattice d, realPoint a in P -> realPoint b in P -> realPoint c in P ->
    Collinear Real ({realPoint a, realPoint b, realPoint c} : Set (Space d)) ->
    a = b \/ a = c \/ b = c
def quasiLinearLengths {d : Nat} (P : Set (Space d)) : Prop :=
  exists k : Nat, 0 < k /\ exists slope intercept : Nat -> Rat,
    forall t : Nat, 0 < t -> exists n : Nat, lengthIs (dilate t P) n /\
      (n : Rat) = slope (t % k) * t + intercept (t % k)
def verticesT : Fin 4 -> Lattice 3 :=
  ![![1, 0, 1], ![3, 3, 1], ![2, 0, 3], ![0, 3, 3]]
def verticesU : Fin 4 -> Lattice 3 :=
  ![![0, 0, 0], ![1, 0, 0], ![0, 0, 1], ![13, 18, 1]]
def problemThreeWitness {d : Nat} (v : Fin (d + 1) -> Lattice d) : Prop :=
  fullSimplex v /\ lengthIs (simplex v) 2 /\ diameterOne (simplex v)
def problemTwoWitness {d : Nat} (v : Fin (d + 1) -> Lattice d) : Prop :=
  fullSimplex v /\ lengthIs (simplex v) 1 /\
    (forall t : Nat, 2 <= t -> lengthIs (dilate t (simplex v)) (3 * t)) /\
    not quasiLinearLengths (simplex v)
def TetrahedronDiameter : Prop := problemThreeWitness verticesT
def TetrahedronDilations : Prop := problemTwoWitness verticesU
def EveryDimensionDiameter : Prop :=
  forall d : Nat, 3 <= d -> exists v : Fin (d + 1) -> Lattice d, problemThreeWitness v
def EveryDimensionDilations : Prop :=
  forall d : Nat, 3 <= d -> exists v : Fin (d + 1) -> Lattice d, problemTwoWitness v
\end{verbatim}}
\noindent Here \texttt{o} stands for composition $\circ$, \texttt{smul} for
scalar multiplication, \verb|''| is the image of a set, \texttt{sum} a
finite sum, \verb|!=| stands for $\ne$, \verb|/\|, \verb|\/|,
\texttt{forall}, \texttt{exists}, \texttt{in}, \verb|<=| and \texttt{not}
stand for $\wedge,\vee,\forall,\exists,\in,\le,\neg$, \verb|![...]| is a vector
literal, and \verb|t % k| is the remainder of $t$ modulo $k$. A point of
\texttt{Lattice d} is a point of $\Z^d$ and \texttt{realPoint} views it in
$\R^d$. \texttt{polytope V} and \texttt{simplex v} are real convex hulls of
lattice points; \texttt{fullSimplex v} says that the $d+1$ vertices are
affinely independent, so the simplex is $d$-dimensional; \texttt{dilate t P}
is $tP$. \texttt{packs P k} says that $P$ contains the Minkowski sum of $k$
lattice polytopes $\conv(V_i)$, each with at least two lattice points, that
is, of positive dimension: every choice of one real point in each factor sums
into $P$. \texttt{lengthIs P n} says that $L(P)=n$: some packing with $n$
summands exists and every packing has at most $n$. \texttt{diameterOne P}
says that $\ell(P)=1$. \texttt{quasiLinearLengths P} says that, for some
period $k>0$ and rational slopes and intercepts depending on $t\bmod k$,
$L(tP)$ exists and equals the corresponding linear function for every
$t\ge1$; this is Definition~1.3 of~\cite{ss}. The four targets are
Theorem~\ref{thm:diam} for $T$, Theorem~\ref{thm:dil} for $U$, and the two
``more generally'' statements.

The final reviewer (see the last section) checked the following points.
(i)~\texttt{packs} is literal containment of the Minkowski sum, as in
Definition~1.1 of~\cite{ss}; each factor sits at its own lattice position,
and translations by non-lattice vectors are not allowed. (ii)~Since
\texttt{lengthIs} asks for an attained maximum, the value is unique (Lean
lemma \texttt{length\_\allowbreak unique}), and the case $k=0$ of
\texttt{packs} plays no role because all values claimed are at least~$1$.
(iii)~\texttt{problemTwoWitness} fixes $L(U)=1$ and $L(tU)=3t$ for $t\ge2$ as
separate conjuncts, so \texttt{not quasiLinearLengths} cannot hold merely
because some maximum fails to exist. The negation is stated for $t\ge1$; a
quasi-linear function on $\N$ including $0$ would restrict to one on $t\ge1$,
and real instead of rational coefficients would make no difference, so the
Lean statement is at least as strong as the negation of the source's notion.
(iv)~\texttt{fullSimplex} makes every witness a genuine $d$-dimensional
simplex in $\R^d$, not a lower-dimensional one padded with zeros.

\paragraph{Review and freezing.}
\texttt{Challenge.lean} was reviewed against~\cite{ss} by a Codex agent that
had written neither the statement nor the proofs, and frozen, by recording its
SHA-256 hash, on 20 November 2026 at 18:06~UTC. The development logs record no
compilation of a proof file before that time (the only earlier compilation is
of \texttt{Challenge.lean} itself, at 18:04~UTC), and the file was not changed
afterwards.

\paragraph{The proofs.}
The two targets for every $d\ge3$ are proved in \texttt{Proof.lean} by
induction on $d$ (\texttt{Nat.le\_\allowbreak induction}), starting from
\texttt{Proof.TetrahedronDiameter} (file \texttt{TTetrahedron.lean}) and
\texttt{Proof.TetrahedronDilations} (file \texttt{UTetrahedron.lean}) and
applying \texttt{pyramid\_\allowbreak length} with $t=1$ and with every
$t\ge2$. In the Lean files the new coordinate of the pyramid is the first one
and the apex is $(1,0,\dots,0)$ (\texttt{pyrVertices}).
Table~\ref{tab:lean} lists the Lean counterparts of the steps of the written
proof (names in the namespace \texttt{TaskP.GeomSimplex.Proof}).

\begin{table}[ht]
\centering\footnotesize
\renewcommand{\arraystretch}{1.08}
\begin{tabular}{@{}>{\raggedright\arraybackslash}p{0.42\textwidth}>{\raggedright\arraybackslash}p{0.54\textwidth}@{}}
\hline
Step of the written proof & Lean (file) \\
\hline
Lemma~\ref{lem:corners} &
\texttt{pack\_endpoints}, \texttt{pack\_of\_endpoints} (\texttt{Packing});
lattice translation: \texttt{translated\_pack} (\texttt{PyramidCorners}) \\
Uniqueness of $L$ & \texttt{length\_unique} (\texttt{Packing}) \\
Lemma~\ref{lem:Tpoints}: barycentric description and dimension of $T$ &
\texttt{tcoords}, \texttt{simplexT\_facets}, \texttt{simplexT\_int},
\texttt{fullT} (\texttt{TCoordinates}) \\
Lemma~\ref{lem:Tpoints}: the eight lattice points &
\texttt{tpoints}, \texttt{points\_of\_T} (\texttt{TCoordinates}) \\
Proposition~\ref{prop:T}, $\ell(T)=1$ (in Lean: a $2\times2$ minor is
nonzero for every triple of distinct points, checked by \texttt{decide}) &
\texttt{diameterT}, from the private \texttt{collinear\_minors} and
\texttt{diameter\_table} (\texttt{TDiameter}) \\
Proposition~\ref{prop:T}, $L(T)\ge2$ and $L(T)\le2$ (in Lean: the
certificate with $\kappa$) &
\texttt{lowerT}, \texttt{upperT}, from the private \texttt{no\_cube} and
\texttt{noCubeTable} (\texttt{TLength}) \\
Equation~\eqref{eq:tU} and dimension of $U$ &
\texttt{uform}, \texttt{budget}, \texttt{simplexU\_facets},
\texttt{dilateU\_facets}, \texttt{dilateU\_int}, \texttt{fullU}
(\texttt{UCoordinates}) \\
Lemma~\ref{lem:Upoints} &
\texttt{length\_oneU}, from the private \texttt{only\_vertices} and
\texttt{no\_parallelogram} (\texttt{UTetrahedron}) \\
Lemma~\ref{lem:cost} &
\texttt{cost\_nonzero}, from the private \texttt{edge\_cost} (\texttt{UUpper}) \\
Proposition~\ref{prop:U}, upper bound & \texttt{upperU} (\texttt{UUpper}) \\
Proposition~\ref{prop:U}, lower bound &
\texttt{lowerU}, from the private \texttt{six}, \texttt{nine},
\texttt{witness\_add}, \texttt{witness\_all} (\texttt{ULower}) \\
Proposition~\ref{prop:U}, not quasi-linear &
\texttt{not\_quasiLinear} (\texttt{Packing}) \\
Lemma~\ref{lem:pyr} &
\texttt{pyrVertices}, \texttt{full\_pyramid}, \texttt{mass\_pyramid},
\texttt{mem\_dilate\_mass} (\texttt{PyramidGeometry});
\texttt{pyramid\_base\_mem}, \texttt{pyramid\_lattice\_cases}
(\texttt{PyramidDiameter}) \\
Proposition~\ref{prop:pyrdiam} &
\texttt{diameter\_pyramid} (\texttt{PyramidDiameter}) \\
Proposition~\ref{prop:pyrlen} &
\texttt{pyramid\_length}, from the private \texttt{lift\_pack} and
\texttt{descend\_pack} (\texttt{PyramidTransport}) and
\texttt{replace\_rising} (\texttt{PyramidCorners}) \\
Induction on the dimension & \texttt{Proof.EveryDimensionDiameter},
\texttt{Proof.EveryDimensionDilations}, private \texttt{one\_level}
(\texttt{Proof}) \\
\hline
\end{tabular}
\caption{Lean counterparts of the steps of the proof.}
\label{tab:lean}
\end{table}

\noindent The test file \texttt{CheckSimplex.lean} restates the four targets
with the frozen types, as four theorems named \texttt{check\_} followed by
the name of the target. The axioms of these four
theorems and of the four \texttt{Proof} theorems, printed by
\texttt{CheckSimplex.lean}, are exactly \texttt{propext},
\texttt{Classical.choice} and \texttt{Quot.sound}. No file uses
\texttt{sorry}, \texttt{native\_\allowbreak decide},
\texttt{bv\_\allowbreak decide} or an \texttt{axiom} declaration. The finite
tables (the $8^3$ triples for $\ell(T)$, the $8^4$ quadruples for $L(T)\le2$
and the $4^4$ quadruples for $L(U)=1$) and the witness inequalities for $2U$
and $3U$ are checked by \texttt{decide}, that is, by the kernel; two of these
proofs raise the \texttt{maxHeartbeats} limit (and one the recursion depth).

\paragraph{Checks.}
All $16$ files, that is, the $13$ proof modules together with the
statement file, \texttt{Proof.lean} and the test file, were compiled module by
module in two new empty directories on the author's Mac, with the pinned
toolchain and with warnings treated as errors, on 20 November 2026
(18:25--18:31~UTC), and every module was replayed with \texttt{leanchecker};
each directory also ran one build of the whole package. All $66$ commands
($33$ per directory) ended with exit code~$0$, the \texttt{leanchecker} runs
printed nothing, and the two directories produced identical \texttt{.olean}
files. The final reviewer (see the next section) compiled the $16$ files again,
together with a short audit file of its own, in a third new directory on the
author's Mac on 20 November 2026 (20:57--21:00~UTC), with warnings treated as
errors, and replayed every module with \texttt{leanchecker}; all runs ended
with exit code~$0$ and the \texttt{leanchecker} runs printed nothing. Mathlib
and its dependencies were taken from a shared prebuilt cache at the pinned
commits and were not rechecked; the \texttt{leanchecker} replays cover the
$16$ project modules. The receipts are in the folder
\texttt{certificates/minkowski-length-simplices} of the repository. In the
repository the files have module names under
\texttt{Research.TaskP.GeomSimplex} and differ from the checked files only in
their \texttt{import} lines (\texttt{Challenge.lean} is unchanged).
The repository versions were made by a script that changed only the local \texttt{import} lines, and on 20 November 2026 (21:04--21:16~UTC) they were compiled and replayed with \texttt{leanchecker} again, module by module, in two new empty directories on the author's Mac, one after the other, together with the files of two other papers of version~\repoversion{} ($48$ modules in all); there were no warnings, and the two directories produced identical \texttt{.olean} files and the same axiom lists. In the whole-project rebuild of version~\repoversion{} of the repository (on the author's Windows computer, in two passes on 20 and 21 November 2026), these files were compiled and replayed with \texttt{leanchecker} again in pass~2 (21 November 2026), which built the files not included in pass~1 on top of its results; no module needed a second attempt (details in the repository). The formalization has not been checked with
comparator or with an independent kernel.

\par}

\section*{How this paper was produced}
This work was carried out from 20 to 22 November 2026 (UTC) in a workflow directed by
the author, using OpenAI Codex with the model GPT-6.2-Astra (reasoning level
ultra) and Anthropic Claude Code with the model Claude Opus~5.5. The author
asked for a further search for substantial open problems, and Claude Code
passed this request to Codex at 17:23~UTC (according to the file timestamp of
the request). Between 17:23 and 17:53~UTC, Codex recorded $111$ problems from
the literature and wrote proofs for five candidates; for the problems of this
paper, the lattice points of $T$, the absence of three-segment sums in $T$,
the segments in $2U$ and $3U$ and the residues in Lemma~\ref{lem:cost} were
checked by exact computation, and the proofs were read by another Codex agent
(same model) that had not taken part in them. The author had asked in advance
that the five most important candidates be formalized at once, without
waiting for a choice; this problem was among the five selected at 18:00~UTC.
From then on, Codex agents (same model) checked the reading of the source,
searched again for prior work, wrote the statement file, the Lean proof and
the written proof, and ran the checks. The reading of the source was checked
by a Codex agent that had not taken part in the written proof, and the
statement file was reviewed before it was frozen by a Codex agent that had
written neither the proofs nor the statement. According to the file
timestamps and the development logs, the proof files were created after the
freeze and compiled between 18:06 and 18:24~UTC; they follow the written proofs
prepared in the search for problems. After the checks, an independent Claude
Code agent (model Claude Opus~5.5) that had written none of the files gave the
final review. It read both arXiv versions of the source itself, compared the
statement file with them, read the written proof, the proof files and the raw
logs of the checks, recomputed the lattice points, the possible segment sums
in $T$ and $U$ and the bound of Lemma~\ref{lem:cost} in Python, and rebuilt
the formalization as described above. Its decision was APPROVE\_FORMAL, with
five minor remarks on the written account (two optional); they were taken
into account in this paper, except the one about the duplicate-check file,
which concerns an internal record. Claude Code (Claude Opus~5.5) checked the sources and wrote this
text from Codex's manuscript and the Lean files. An independent Claude Code agent (model Claude Opus~5.5) that had written none of the files reviewed this paper, and another reviewed the release; their corrections were applied. Throughout,
``we'' refers to this workflow; the reading and searching described in the
paper were done by the agents. Codex agents searched for prior work on
20 November 2026, in the search for problems and again before the statement
was frozen, and the final reviewer read the closest later works it found,
including~\cite{bs,mss}; they found no answer to Problems~2 or~3
of~\cite{ss} (details in the repository). Claude Code searched again on 21 November 2026 (17:54--18:21~UTC; the arXiv, Semantic Scholar, Crossref, DataCite, Zenodo, MathDB, GitHub and repositories of AI-assisted mathematics, including the works citing~\cite{mss} and its published code) and found no new related work; that day the citation index OpenAlex refused to list the works citing~\cite{ss} (rate limit). Claude Code searched again on 22 November 2026 (05:39--05:51~UTC), when the OpenAlex lists of citing works refused on 21 November were retrieved, and found nothing new. This reports the
coverage of the searches, not a guarantee that they were complete. The author
is not a professional mathematician. The author has read this paper and, using a side-by-side table, compared its Lean statements with the main theorems: Theorems~\ref{thm:diam} and~\ref{thm:dil}. No human mathematician has
reviewed this paper. The correctness of the results rests on the Lean proofs
and on the agreement of the Lean statements with the problems of~\cite{ss};
the written proofs are given for the reader. AI systems are not authors of
this paper; the author takes full responsibility for its content.

\medskip
\noindent Thanks are due to Ivan Soprunov and Jenya Soprunova for the
questions.

{\small
\begin{thebibliography}{9}
\setlength{\itemsep}{1pt}
\bibitem{bgsw} O.~Beckwith, M.~Grimm, J.~Soprunova and B.~Weaver, Minkowski
length of 3D lattice polytopes, \emph{Discrete Comput. Geom.} 48 (2012),
no.~4, 1137--1158, \url{https://doi.org/10.1007/s00454-012-9433-5}.
\bibitem{bs} M.~Blanco and F.~Santos, Enumeration of lattice 3-polytopes by
their number of lattice points, \emph{Discrete Comput. Geom.} 60 (2018),
no.~3, 756--800, \url{https://doi.org/10.1007/s00454-017-9932-5};
arXiv:1601.02577v3.
\bibitem{repo} H.~Dong, \emph{ai4math-results}, repository,
\url{https://github.com/Horace-Maxwell/ai4math-results}, version~\repoversion{}
(2026); archived at Zenodo under the concept DOI
\url{https://doi.org/10.5281/zenodo.22970254}.
\bibitem{mathlib}
The mathlib Community, The Lean mathematical library, in \emph{Proc.\ 9th ACM
SIGPLAN Int.\ Conf.\ Certified Programs and Proofs} (CPP '20), ACM, 2020,
\url{https://doi.org/10.1145/3372885.3373824}.
\bibitem{mss} K.~Meyer, I.~Soprunov and J.~Soprunova, $\mathbb F_q$-zeros of
sparse trivariate polynomials and toric 3-fold codes, \emph{SIAM J. Appl.
Algebra Geom.} 6 (2022), no.~3, 432--467,
\url{https://doi.org/10.1137/21M1436890}; arXiv:2105.10071v2.
\bibitem{ss} I.~Soprunov and J.~Soprunova, Eventual quasi-linearity of the
Minkowski length, \emph{European J. Combin.} 58 (2016), 107--117,
\url{https://doi.org/10.1016/j.ejc.2016.05.009}; arXiv:1412.4404v2 (2016).
\bibitem{white} G.~K. White, Lattice tetrahedra, \emph{Canad. J. Math.} 16
(1964), 389--396, \url{https://doi.org/10.4153/CJM-1964-040-2}.
\end{thebibliography}}
\end{document}
